\documentclass{article}
\usepackage[12pt]{extsizes}
\usepackage[a4paper, margin=0.5cm]{geometry}
\usepackage{indentfirst}
\usepackage[T1, T2A]{fontenc}
\usepackage[english, russian]{babel}
\usepackage{amsmath, amssymb, fleqn, cancel, graphicx, wrapfig, tabbing, color, cases, lipsum}

\usepackage{amstext} % for \text macro
\usepackage{array}   % for \newcolumntype macro
\newcolumntype{L}{>{$}l<{$}} % math-mode version of "l" column type

\DeclareMathOperator{\rank}{rank}
\makeatletter
\newenvironment{sqcases}{%
  \matrix@check\sqcases\env@sqcases
}{%
  \endarray\right.%
}
\def\env@sqcases{%
  \let\@ifnextchar\new@ifnextchar
  \left\lbrack
  \def\arraystretch{1.2}%
  \array{@{}l@{\quad}l@{}}%
}
\makeatother

\DeclareRobustCommand{\divby}{\mathrel{\vbox{\baselineskip.65ex\lineskiplimit0pt\hbox{.}\hbox{.}\hbox{.}}}}

\setlength{\parindent}{0pt}

\newenvironment{rcases}
    {\left.\begin{aligned}}
    {\end{aligned}\right\rbrace}

\usepackage{hyperref}
\hypersetup{
    colorlinks,
    citecolor=black,
    filecolor=black,
    linkcolor=black,
    urlcolor=black
}


% Примечание* --- информация, добавленная от себя, но не написанная преподавателем на доске
% 99.9% Материала в этом и других документах является информацией, набранной с доски, и, если меня не было на занятии, набранной с конспектов одногруппников.
% Поэтому в моих документах редко где может встретиться то, что преподаватель сказал, но не написал (или одногруппники что-то недописали).


\begin{document}

\begin{titlepage}
    \title{Линейная алгебра и Аналитическая геометрия}
    \author{Бабицкий Иван | М8О-111БВ-24 \\ 02.03.02 "Фундаментальная информатика и информационные технологии"}
    \date{2024/2025 учебный год, 1-й семестр}
    \maketitle
    \begin{center}
        
\includegraphics[scale=0.2]{images/mai.png}
    \end{center}
\end{titlepage}

\tableofcontents



%--------------------------------------------------
\newpage
\paragraph{ПЗ-9 | 2024.11.11 (3)}
\part{Произведения геометрических векторов}

\section{Теория}

\subsection{Скалярное}

\subsubsection{Определение}
\begin{gather*}
    \text{Число } = (\overline{a}, \overline{b}) \equiv \overline{a} \cdot \overline{b} \overset{\Delta}{=}
    |\overline{a}| \cdot \underbrace{|\overline{b}| \cdot \cos \varphi}_{\text{Пр}_{\overline{a}} \overline{b}}
    ~~~~~ \varphi = \widehat{\overline{a}, \overline{b}} \in \left[ 0; \pi \right] 
\end{gather*}

\subsubsection{Геометрический смысл}
(рисунок)

\subsubsection{Вычисление}
\begin{align*}
    \begin{rcases}
        \overline{a} = a_x \overline{i} + a_y \overline{j} + a_z \overline{k} = \{a_x, a_y, a_z\}_{(\overline{i}, \overline{j}, \overline{k})} \\
        \overline{b} = b_x \overline{i} + b_y \overline{j} + b_z \overline{k} = \{b_x, b_y, b_z\}_{(\overline{i}, \overline{j}, \overline{k})}
    \end{rcases} \Rightarrow
    \overline{a} \cdot \overline{b} = \boxed{a_x b_x + a_y b_y + a_z b_z}
\end{align*}

\subsubsection{Свойства}
\begin{itemize}
    \item 
    \[
        \overline{a} \cdot \overline{b} = \overline{b} \cdot \overline{a}
    \]
    \item 
    \[
    \overline{a} \cdot \overline{b} = 0, 
    \text{ если } \overline{a} \neq \overline{0}, 
    \overline{b} \neq \overline{0}, 
    \underbrace{\overline{a} \bot \overline{b}}_{\text{ортогональные}}
    \]
    \item 
    \[
    \overline{a} \cdot \overline{a} = |\overline{a}|^2
    \]
    \item 
    \[
    (\alpha_1 \overline{a}_1 + \alpha_2 \overline{a}_2) \cdot \overline{b} =
    \alpha_1 \underbrace{\overline{a}_1 \overline{b}}_{= \overline{b} \overline{a}_1} + \alpha_2 \overline{a}_2 \overline{b}
    \]
    \item 
    \[
    \overline{a} \cdot \overline{b} = |\overline{a}| \cdot \text{Пр}_{\overline{a}} \overline{b}
    \]
\end{itemize}

\subsubsection{Применение}
\begin{itemize}
    \item 
    \[
    |\overline{a}| = \sqrt{\overline{a} \cdot \overline{a}} = \sqrt{a_x^2 + a_y^2 + a_z^2}
    \]
    \item 
    \[
    \text{Пр}_{\overline{a}} \overline{b} =\frac{\overline{a} \cdot \overline{b}}{|\overline{a}|}
    \]
    \item 
    \[
    \cos \varphi = \frac{\overline{a} \cdot \overline{b}}{|\overline{a}| \cdot |\overline{b}|}
    \]
\end{itemize}


\par\noindent\rule{\textwidth}{1pt}
\subsection{Векторное}

\subsubsection{Определение}
\[
    \text{Вектор } = \left[ \overline{a}, \overline{b} \right] \equiv \overline{a} \times \overline{b} = \overline{n}:
\]
\begin{itemize}
    \item $\overline{n} \bot \overline{a}, \overline{n} \bot \overline{b}$
    \item Тройка $\overline{a}, \overline{b}, \overline{n}$ --- правая ($\overline{n}$ --- по правилу буравчика)
    \item $|\overline{n}| = |\overline{a}| \cdot |\overline{b}| \cdot \sin \varphi = S_{\#}$
\end{itemize}


\subsubsection{Геометрический смысл}
\includegraphics[scale=0.15]{images/2024-11-11/01.jpg}

\subsubsection{Вычисление (через декартовы прямые координаты)}
\begin{gather*}
    \overline{n} = \overline{a} \times \overline{b} =
    \begin{vmatrix}
        \overline{i} & \overline{j} & \overline{k} \\
        a_x          & a_y          & a_z          \\
        b_x          & b_y          & b_z
    \end{vmatrix} =
    A_{11} \overline{i} + A_{12} \overline{j} + A_{13} \overline{k} =
    \{\underbrace{A_{11}}_{n_x}, \underbrace{A_{12}}_{n_y}, \underbrace{A_{13}}_{n_z}\}
\end{gather*}


\subsubsection{Свойства}
\begin{itemize}
    \item 
    \[
    \overline{a} \times \overline{b} = -\overline{b} \times \overline{a}
    \]
    \item 
    \begin{align*} 
        \begin{rcases}
            \overline{a} \times \overline{b} = \overline{0} \\
            \overline{a} \neq \overline{0}, \overline{b} \neq \overline{0}    
        \end{rcases} \Leftrightarrow 
        \overline{a} \parallel \overline{b} 
    \end{align*}
    \item 
    \[
    \overline{a} \times \overline{a} = \overline{0}
    \]
    \item 
    \[
    \left( \alpha_1 \overline{a}_1 + \alpha_2 \overline{a}_2 \right) \times \overline{b} =
    \alpha_1 \underbrace{\overline{a}_1 \times \overline{b}}_{\neq \overline{b} \times \overline{a}_1} +
    \alpha_2 \overline{a}_2 \times \overline{b}
    \]
\end{itemize}

\subsubsection{Применение}
\begin{itemize}
    \item 
    \[
    S_{\#}^{\overline{a}, \overline{b}} = \left| \overline{a} \times \overline{b} \right|
    \]
    \item 
    \[
    \overline{n} = \overline{a} \times \overline{b} \text{ --- нормаль к плоскости } P \parallel \overline{a}, \overline{b}
    \]
\end{itemize}


\par\noindent\rule{\textwidth}{1pt}
\subsection{Смешанное (векторно-скалярное)}

\subsubsection{Определение}
\begin{gather*}
    \text{Число } = (\overline{a}, \overline{b}, \overline{c}) \overset{\Delta}{=}
    (\overbrace{\overline{a} \underset{(1)}{\times} \overline{b}}^{\overline{n}}) \underset{(2)}{\cdot} \overline{c}
\end{gather*}


\subsubsection{Геометрический смысл}
(рисунок)

\subsubsection{Вычисление}
\begin{gather*}
    (\overline{a}, \overline{b}, \overline{c}) =
    \begin{vmatrix}
        a_x & a_y & a_z \\
        b_x & b_y & b_z \\
        c_x & c_y & c_z
    \end{vmatrix}
\end{gather*}
\paragraph{Проверка по $\overline{n} = \overline{a} \times \overline{b}$}
\[
    \overline{n} \cdot \overline{c} \overset{?}{=} (\overline{a}, \overline{b}, \overline{c})
\]


\subsubsection{Свойства}
\begin{itemize}
    \item Перестановки круговые:
    \[
    (\overline{a}, \overline{b}, \overline{c}) =
    (\overline{b}, \overline{c}, \overline{a}) =
    (\overline{c}, \overline{a}, \overline{b}) =
    (\overline{a}, \overline{b}, \overline{c})
    \]
    \item Перестановки соседей:
    \[
     (\overline{a}, \overline{b}, \overline{c}) =
    -(\overline{b}, \overline{a}, \overline{c})
    \]
    \[
     (\overline{a}, \overline{b}, \overline{c}) =
    -(\overline{a}, \overline{c}, \overline{b})
    \]
    \item 
    \[
    \left( \left( \alpha_1 \overline{a}_1 + \alpha_2 \overline{a}_2 \right), \overline{b}, \overline{c} \right) =
    \alpha_1 \left( \overline{a}_1, \overline{b}, \overline{c} \right) +
    \alpha_2 \left( \overline{a}_2, \overline{b}, \overline{c} \right)
    \]
    \item Компланарная тройка:
    \begin{align*}
        \begin{rcases}
            (\overline{a}, \overline{b}, \overline{c}) = 0 \\
            \overline{a} \neq \overline{0}, \overline{b} \neq \overline{0}, \overline{c} \neq \overline{0} 
        \end{rcases} \Leftrightarrow
        \overline{a}, \overline{b}, \overline{c} \text{ --- тройка компланарна}
    \end{align*}
\end{itemize}

\subsubsection{Применение}
\begin{itemize}
    \item Объём параллелепипеда:
    \[
    V_{\#}^{\overline{a}, \overline{b}, \overline{c}} = \left| \left( \overline{a}, \overline{b}, \overline{c} \right) \right|
    \]
    \item
    \begin{itemize}
        \item $\left( \overline{a}, \overline{b}, \overline{c} \right) > 0 \Leftrightarrow$ тройка правая
        \item $\left( \overline{a}, \overline{b}, \overline{c} \right) < 0 \Leftrightarrow$ тройка левая
    \end{itemize}
\end{itemize}



\par\noindent\rule{\textwidth}{1pt}
\section{Практика}

\subsection{1.108}

\paragraph{Дано}
\begin{itemize}
    \item $\Delta ABC$
    \item $A(1; -1; 2)$
    \item $B(5; -6; 2)$
    \item $C(1; 3; -1)$
\end{itemize}

\paragraph{Найти}
\begin{itemize}
    \item $\angle A = \alpha = \widehat{\overline{AB}, \overline{AC}}$
    \item $S_{ABC}$
    \item Высота $h_B = BD, BD \bot AC$
\end{itemize}

\paragraph{Решение}
\begin{enumerate}
    \item Введём векторы сторон $\Delta ABC$, исходящие из одной вершины:
    $\overline{AB} = \{5 - 1; -6 - (-1); 2 - 2\} = \{4; -5; 0\} = 4\overline{i} - 5\overline{j} + 0\overline{k}$ \\
    $\overline{AC} = \{0; 4; - 3\} = 0\overline{i} + 4\overline{j} - 3\overline{k}$
    \item $\cos \alpha = \frac{\overline{AB} \cdot \overline{AC}}{|\overline{AB}| \cdot |\overline{AC}|} = 
    \frac{4 \cdot 0 - 4 \cdot 5 - 3 \cdot 0}{\sqrt{16 + 25} \cdot \sqrt{16 + 9}} = 
    \frac{-20}{5\sqrt{41}} = -\frac{4}{\sqrt{41}} \in [-1; 1] \\
    \alpha = \arccos \left( \frac{-4}{\sqrt{41}} \right) = \pi - \arccos \left( \frac{4}{\sqrt{41}} \right)$ --- тупой
\end{enumerate}




%--------------------------------------------------
\newpage
\paragraph{ПЗ-10 | 2024.11.18 (3)}
\part{Прямая линия $l$ в плоскости $Oxy$}

\section{Теория}

\subsection{Общее уравнение}
\includegraphics[scale=0.2]{images/2024-11-18/01.jpg}
\begin{gather*}
    \forall M(x, y) \in l: \overline{M_0 M} \bot \overline{N} \Leftrightarrow \overline{N} \cdot \overline{M_0 M} = 0 \\
    \boxed{A(x- x_0) + B(y - y_0) = 0} \Rightarrow \boxed{Ax + By + C = 0} \\
\end{gather*}

\subsubsection{Примечание*}
$\overline{N}(A, B)$ --- вектор нормали прямой $l$.

\subsubsection{Случаи}
\begin{enumerate}
    \item $A = 0 \Rightarrow y = y_{*}$
    \item $B = 0 \Rightarrow x = x_{*}$
    \item $C = 0 \Rightarrow y = kx$
\end{enumerate}
\includegraphics[scale=0.2]{images/2024-11-18/02.jpg}


\par\noindent\rule{\textwidth}{1pt}
\subsection{Параметрические и каноническое уравнения $l \subset Oxy$}
\includegraphics[scale=0.2]{images/2024-11-18/03.jpg}
\begin{gather*}
    \forall M(x, y) \in l: \overline{M_0 M} \parallel \overline{q} \Leftrightarrow \exists t \in \mathbb{R}: \overline{M_0 M} = t\overline{q} \Leftrightarrow
    \begin{cases}
        x - x_0 = mt \\
        y - y_0 = nt
    \end{cases} t \in \mathbb{R} \Rightarrow \\
    \Rightarrow t = \frac{x - x_0}{m} = \frac{y - y_0}{n} \Rightarrow
    n(x - x_0) = m(y - y_0) \Rightarrow \boxed{n(x - x_0) - m(y - y_0) = 0}
\end{gather*}

\subsubsection{Примечание*}
$\overline{q}(-B, A)$ --- направляющий вектор прямой $l$.

\subsubsection{Случай 1}
\[
    m = 0
\]
\begin{gather*}
    \frac{x - 1}{0} = \frac{y + 2}{3} \\
    \begin{cases}
        x - 1 = 0t \rightarrow x = 1 \\
        y + 2 = 3t \rightarrow y \in \mathbb{R}
    \end{cases}
\end{gather*}
\includegraphics[scale=0.15]{images/2024-11-18/04.jpg}

\subsubsection{Случай 2}
\[
    n = 0
\]
\begin{gather*}
    \frac{x - 2}{4} = \frac{y + 5}{0} \\
    \begin{cases}
        x - 2 = 4t \rightarrow x \in \mathbb{R} \\
        y + 5 = 0t \rightarrow y = -5
    \end{cases}
\end{gather*}
\includegraphics[scale=0.15]{images/2024-11-18/05.jpg}


\par\noindent\rule{\textwidth}{1pt}
\subsection{Через 2 точки}
\includegraphics[scale=0.2]{images/2024-11-18/06.jpg}
\begin{align*}
    \begin{rcases}
        M_1(x_1, y_1) \in l \\
        M_2(x_2, y_2) \in l
    \end{rcases} \Rightarrow
    \boxed{\frac{x - x_1}{x_2 - x_1} = \frac{y - y_1}{y_2 - y_1}}
\end{align*}


\par\noindent\rule{\textwidth}{1pt}
\subsection{Расстояние от точки $D$ до прямой $l$}
\begin{gather*}
    l: Ax + By + C = 0, ~~~ D(x_D, y_D) \\
    \rho(D(x_D, y_D); l: Ax + By + C = 0) = \frac{|A_{x_D} + B_{y_D} + C|}{\sqrt{A^2 + B^2}}
\end{gather*}
\includegraphics[scale=0.15]{images/2024-11-18/07.jpg}


\par\noindent\rule{\textwidth}{2pt}
\section{Практика}

\subsection{1.150}

\paragraph{Дано}
\begin{itemize}
    \item $\Delta ABC$
    \item $A(1; 2)$
    \item $B(2; -2)$
    \item $C(6; 1)$
\end{itemize}

\paragraph{Найти}
\begin{enumerate}
    \item Уравнение стороны $(AB)$
    \item Уравнение высоты $(CD)$
    \item Длина $h_c$ высоты $CD$ 
    \item Уравнение медианы $(BM)$
    \item Точка пересечения $k = (CD) \cap (BM)$
    \item Угол между высотой $(CD)$ и медианой $(BM)$
\end{enumerate}

\paragraph{Примечание} 
~\\ Проверить найденные результаты по рисунку.

\paragraph{Пункт 1 | Решение}
\begin{align*}
    \begin{rcases}
        M_1 = A(1; 2)  \in (AB) \\
        M_2 = B(2; -2) \in (AB)
    \end{rcases} \Rightarrow
    \text{каноническое уравнение: } \frac{x - x_1}{x_2 - x_1} = \frac{y - y_1}{y_2 - y_1}
\end{align*}
\begin{gather*}
    \frac{x - 1}{2 - 1} = \frac{y - 2}{-2 - 2} ~~\rightarrow~~
    \frac{x - 1}{1} = \frac{y - 2}{-4} ~~\rightarrow~~
    -4(x - 1) = y - 2 ~~\rightarrow~~
    y - 2 + 4(x - 1) = 0 \\
    y - 2 + 4x - 4 = 0 ~~\rightarrow~~
    \boxed{4x + y - 6 = 0} \text{ ---  общее уравнение}
\end{gather*}

\paragraph{Пункт 1 | Проверка}
Проверка точек пересечения с осями
\begin{gather*}
    Ox: y = 0 \Rightarrow 4x + 0 - 6 = 0 ~~\rightarrow~~ x = \frac{6}{4} = \frac{3}{2} \text{ --- есть на рисунке} \\
    Oy: x = 0 \Rightarrow 4 \cdot 0 + y - 6 = 0 ~~\rightarrow~~ y = 6 \text{ --- есть на рисунке}
\end{gather*}

\paragraph{Пункт 1 | Ответ}
\[
    (AB): 4x + y - 6 = 0
\]

\paragraph{Пункт 2 | Решение}
\begin{gather*}
    \begin{cases}
        M_0 = C(6; 1) \in (CD) \\
        \frac{\pi}{2}: \overline{AB} \bot \overline{CD} \Rightarrow N_{(CD)} = \overline{AB} = \{2 - 1; -2 - 2\} = \{1; -4\}
    \end{cases} \\
    \text{Общее уравнение с точкой: } \\
    A(x - x_0) + B(y - y_0) = 0 ~~\rightarrow~~
    1(x - 6) - 4(y - 1) = 0 ~~\rightarrow~~
    \boxed{x - 4y - 2 = 0}
\end{gather*}

\paragraph{Пункт 2 | Проверка}
\begin{gather*}
    Oy: x = 0 \Rightarrow y = -\frac{1}{2} \text{ --- есть на рисунке} \\
    Ox: y = 0 \Rightarrow x = 2 \text{ --- есть на рисунке}
\end{gather*}

\paragraph{Пункт 2 | Ответ}
\[
    (CD): x - 4y - 2 = 0
\]

\paragraph{Пункт 3 | Решение}
\begin{align*}
    \begin{rcases}
        \text{точка } C(6; 1) = D \\
        \text{уравнение } (AB): 4x + y - 6 = 0
    \end{rcases} \Rightarrow \\
\end{align*}
\begin{gather*}
    \Rightarrow h_c = \rho(C, (AB)) = 
    \frac{\left| Ax_d + By_d + C \right|}{\sqrt{A^2 + B^2}} = 
    \frac{\left| 4 \cdot 6 + 1 \cdot 1 - 6 \right|}{\sqrt{4^2 + 1^2}} = 
    \boxed{\frac{19}{\sqrt{17}}} \approx 
    4.6 \text{ --- есть на рисунке}
\end{gather*}

\paragraph{Пункт 3 | Ответ}
\[
    h_c = \frac{19}{\sqrt{17}}
\]

\paragraph{Пункт 4 | Решение}
\begin{align*}
    \begin{rcases}
        \text{точка } B(2; - 2) \in BM \\
        \text{точка } M\left( \frac{6 + 1}{2}; \frac{1 + 2}{2} \right) \Rightarrow \text{точка } M(3,5; 1,5)
    \end{rcases} \Rightarrow
\end{align*}
\begin{gather*}
    \Rightarrow \frac{x - x_1}{x_2 - x_1} = \frac{y - y_1}{y_2 - y_1} ~~\rightarrow~~ 
    \frac{x - 2}{3,5 - 2} = \frac{y + 2}{1,5 + 2} ~~\rightarrow~~
    \frac{x - 2}{1,5} = \frac{y + 2}{3,5} ~~\rightarrow~~
    \frac{2x - 4}{3} = \frac{2y + 4}{7} \\
    14x - 28 = 6y + 12 ~~\rightarrow~~
    14x - 6y - 40 = 0 ~~\rightarrow~~
    \boxed{7x - 3y - 20 = 0}
\end{gather*}

\paragraph{Пункт 4 | Проверка}
\begin{gather*}
    Oy: x = 0 \Rightarrow y = -\frac{20}{3} \approx -7 \text{ --- есть на рисунке} \\
    Ox: y = 0 \Rightarrow x = \frac{20}{7} \approx 3 \text{ --- есть на рисунке}
\end{gather*}

\paragraph{Пункт 4 | Ответ}
\[
    (BM): 7x - 3y - 20 = 0
\]

\paragraph{Пункт 5 | Решение и проверка}
\begin{align*}
    k = (CD) \cap (BM):
    \begin{rcases}
        k \in (CD) \\
        k \in (BM)
    \end{rcases} \Rightarrow
\end{align*}
\begin{gather*}
    \Rightarrow K(x, y) - ?
    \begin{cases}
        x - 4y - 2 = 0 \\
        7x - 3y - 20 = 0
    \end{cases} \rightarrow~~
    \begin{cases}
        x - 4y = 2 \\
        7x - 3y = 20
    \end{cases}
\end{gather*}
Правило Крамера:
\begin{gather*}
    A = \begin{pmatrix}
        1 & -4 \\
        7 & -3
    \end{pmatrix}, ~~
    \overline{b} = \begin{pmatrix}
        2  \\
        20
    \end{pmatrix}, ~~
    A \cdot \overline{x} = \overline{b}, ~~
    x_i = \frac{\Delta_i}{\Delta}, ~~
    \Delta = \det A \\
    \Delta = -3 + 25 = 25 \neq 0 \Rightarrow \exists A^{-1}! \\
    \Delta_1 = \begin{vmatrix}
        2  & -4 \\
        20 & -3
    \end{vmatrix} = -6 + 80 = 74 \\
    \Delta_2 = \begin{vmatrix}
        1 &  2 \\
        7 & 20
    \end{vmatrix} = 20 - 14 = 6 \\
    x_1 = \frac{74}{25} = 2.96 \\
    x_2 = \frac{6}{25}  = 0.24  \\
    k(3 - \frac{1}{25}; 0.24) \text{ --- есть на рисунке}
\end{gather*}

\paragraph{Пункт 5 | Ответ}
\[
    k(2.96; 0.24)
\]

\paragraph{Пункт 6 | Решение и проверка}
\begin{align*}
    \begin{rcases}
        (CD): x - 4y - 2 = 0   \Rightarrow \overline{N_1} = \{1; -4\} \bot (CD) \\
        (BM): 7x - 3y - 20 = 0 \Rightarrow \overline{N_2} = \{7; -3\} \bot (BM)
    \end{rcases} \Rightarrow
\end{align*}
\begin{gather*}
    \Rightarrow \cos \alpha = |\cos \varphi| = \left| \frac{\overline{N_1} \cdot \overline{N_2}}{|\overline{N_1}| \cdot |\overline{N_2}|} \right| =
    \frac{| \{1; -4\} \cdot \{7; -3\} |}{| \{1; -4\} | \cdot | \{7; -3\} | } =
    \frac{|1 \cdot 7 + 4 \cdot 3|}{\sqrt{1^2 + 4^2} \cdot \sqrt{7^2 + 3^2}} =
    \frac{19}{\sqrt{17} \cdot \sqrt{58}} \in [0; 1] \\
    \alpha = \arccos \frac{19}{\sqrt{17} \cdot \sqrt{58}} = \dots \approx 53^{\circ} \text{ --- совпадает с рисунком}
\end{gather*}

\paragraph{Пункт 6 | Ответ}
\[
    \alpha \approx 53^{\circ}
\]




%--------------------------------------------------
\newpage
\paragraph{ПЗ-11 | 2024.11.25 (3)}
\part{Плоскость $P$ в пространстве $Oxyz$}

\section{Теория}

\subsection{Общее уравнение}
\includegraphics[scale=0.15]{images/2024-11-25/01.jpg}
\begin{align*}
    \begin{rcases}
        \overline{N} = \{A, B, C\} \bot P \\
        M_0(x_0, y_0, z_0) \in P
    \end{rcases} \Rightarrow
    \forall M(x, y, z) \in P: \overline{M_0 M} \bot \overline{N} \Leftrightarrow \overline{N} \cdot \overline{M_0 M} = 0 \Leftrightarrow
\end{align*}
\begin{gather*}
    \Leftrightarrow \boxed{A(x - x_0) + B(y - y_0) + C(z - z_0) = 0} ~~\rightarrow~~ \boxed{Ax + By + Cz + D = 0} \Rightarrow \{A, B, C\} \bot P
\end{gather*}

\subsubsection{Случай 1}
\[
    A = 0 ~~~~~ By + Cz + D = 0: l_{yz}
\]
\includegraphics[scale=0.15]{images/2024-11-25/02.jpg}

\subsubsection{Случай 2}
\[
    A = B = 0 ~~~~~ Cz = D ~~~~~ z = z_*
\]
\includegraphics[scale=0.15]{images/2024-11-25/03.jpg}


\par\noindent\rule{\textwidth}{1pt}
\subsection{Компланарное уравнение}
\includegraphics[scale=0.15]{images/2024-11-25/04.jpg}
\begin{align*}
    \begin{rcases}
        M_0(x_0, y_0, z_0) \in P \\
        \overline{p} \parallel P \\
        \overline{q} \parallel P \\
        \overline{p} \nparallel q
    \end{rcases} \Rightarrow
    \forall M(x, y, z) \in P: \overline{M_0 M}, \overline{p}, \overline{q} \text{ --- компланарны} \Leftrightarrow
    (\overline{M_0 M}, \overline{p}, \overline{q}) = 0
\end{align*}
\begin{gather*}
    \begin{vmatrix}
        x - x_0 & y - y_0 & z - z_0 \\
        p_x     & p_y     & p_z     \\
        q_x     & q_y     & q_z
    \end{vmatrix} = 0 \Rightarrow
    (x - x_0)A_{11} + (y - y_0)A_{12} + (z - z_0)A_{13} = 0
\end{gather*}


\par\noindent\rule{\textwidth}{1pt}
\subsection{Через 3 точки}
\includegraphics[scale=0.15]{images/2024-11-25/05.jpg}
\begin{gather*}
    P: \begin{vmatrix}
        x   - x_1 & y   - y_1 & z   - z_1 \\
        x_2 - x_1 & y_2 - y_1 & z_2 - z_1 \\
        x_3 - x_1 & y_3 - y_1 & z_3 - z_1 \\
    \end{vmatrix} = 0
\end{gather*}


\par\noindent\rule{\textwidth}{1pt}
\subsection{Расстояние от точки $M_1$ до плоскости $P$}
\begin{gather*}
    P: Ax + By + Cz + D = 0 \\
    \rho(M_1, P) = \frac{| Ax_1 + By_1 + Cz_1 + D |}{\sqrt{A^2 + B^2 + C^2}}
\end{gather*}


\par\noindent\rule{\textwidth}{1pt}
\subsection{Взаимное расположение плоскостей $P_1$ и $P_2$}
\begin{gather*}
    P_1: A_1x + B_1y + C_1z + D_1 = 0 \\
    P_2: A_2x + B_2y + C_2z + D_2 = 0
\end{gather*}

\subsubsection{Случай 1}
\begin{gather*}
    \boxed{P_1 \parallel P_2 } \Leftrightarrow
    \frac{A_1}{A_2} = \frac{B_1}{B_2} = \frac{C_1}{C_2}:
    \rho(P_1 \parallel P_2) = \rho(M_1 \in P_1, P_2) =
    \frac{| A_2 x_1 + B_2 y_1 + C_2 z_1 + D_2 |}{\sqrt{A^2_2 + B^2_2 + C^2_2}}, \\
    \text{где } x_1, y_1, z_1 ?: \begin{cases}
        A_1 x_1 + B_1 y_1 + C_1 z_1 + D_1 = 0 \\
        \underbrace{x_1 = y_1 = 0}_{z_1 - ?} \text{ или } \underbrace{x_1 = z_1 = 0}_{y_1 - ?} \text{ или } \underbrace{y_1 = z_1 = 0}_{x_1 - ?}
    \end{cases}
\end{gather*}

\subsubsection{Случай 2}
\begin{gather*}
    \boxed{P_1 \nparallel P_2} \Leftrightarrow 
    \frac{A_1}{A_2} \neq \frac{B_1}{B_2} \text{ или } \frac{B_1}{B_2} \neq \frac{C_1}{C_2} \\
    \alpha = \widehat{P_1, P_2} \in \left( 0; \frac{\pi}{2} \right]:
    \cos \alpha = |\cos \varphi| = \frac{| \overline{N_1} \cdot \overline{N_2} |}{| \overline{N_1} | \cdot | \overline{N_2} |}
\end{gather*}



\par\noindent\rule{\textwidth}{2pt}
\section{Практика}

\subsection{}


%--------------------------------------------------
\newpage
\paragraph{ПЗ-12 | 2024.12.02 (3)}
\part{Прямая линия $L$ в пространстве $Oxyz$}

\section{Теория}

\subsection{Общие уравнения}
\begin{gather*}
    L:
    \begin{cases}
        A_1x + B_1y + C_1z + D_1 = 0 \rightarrow \overline{N_1} \\
        A_2x + B_2y + C_2z + D_2 = 0 \rightarrow \overline{N_2}
    \end{cases} \\
    \frac{A_1}{A_2} \neq \frac{B_1}{B_2} \text{ или } \frac{B_1}{B_2} \neq \frac{C_1}{C_2} \\
\end{gather*}


\subsection{Параметрические и каноническое уравнения}
\begin{gather*}
    \begin{cases}
        x - x_0 = mt \\
        y - y_0 = nt \\
        z - z_0 = pt
    \end{cases} t \in \mathbb{R} 
    \Leftrightarrow t = \frac{x - x_0}{m} = \frac{y - y_0}{n} = \frac{z - z_0}{p}
\end{gather*}


\subsection{Расстояние от точки $M_1$ до прямой линии $L$}
\begin{gather*}
    L: \frac{x - x_0}{m} = \frac{y - y_0}{n} = \frac{z - z_0}{p} \\
    \rho(M_1, L) = \frac{\left| \overline{M_0 M_1} \times \overline{q} \right|}{\left| \overline{q} \right|}
\end{gather*}


\subsection{Взаимное расположение $L_1$ и $L_2$}
\begin{gather*}
    L_1: \frac{x - x_1}{m} = \frac{y - y_1}{n} = \frac{z - z_1}{p} \\
    L_2: \frac{x - x_2}{m} = \frac{y - y_2}{n} = \frac{z - z_2}{p} \\
    \text{Из числителей: }   M_1 \in L_1, M_2 \in L_2 \\
    \text{Из знаменателей: } \overline{q_1} \parallel L_1, \overline{q_2} \parallel L_2 \\
\end{gather*}

\subsubsection{Случай 1}
\[
    L_1 \parallel L_2 \Leftrightarrow \frac{m_1}{m_2} = \frac{n_1}{n_2} = \frac{p_1}{p_2}
\]
\begin{gather*}
    \rho(L_1 \parallel L_2) = 
    \rho(M_1 \in L_1, L_2) = 
    \frac{\left| \overline{M_2 M_1} \times \overline{q_2} \right|}{\left| \overline{q_2} \right|} =
    \begin{cases}
        \neq 0 \text{ --- не совпадают} \\
           = 0 \text{ --- совпадают}
    \end{cases}
\end{gather*}

\subsubsection{Случай 2}
\[
    L_1 \nparallel L_2 \Leftrightarrow \frac{m_1}{m_2} \neq \frac{n_1}{n_2} \text{ или } \frac{n_1}{n_2} \neq \frac{p_1}{p_2}
\]
\begin{gather*}
    \rho(L_1 \nparallel L_2) =
    \frac{\left| \left( \overline{M_2 M_1}, \overline{q_1}, \overline{q_2} \right) \right|}{\left| \overline{q_1} \times \overline{q_2} \right|} =
    \begin{cases}
        \neq 0 \text{ --- скрещиваются} \\
           = 0 \text{ --- пересекаются}
    \end{cases}
\end{gather*}
\begin{gather*}
    \alpha = \widehat{L_1, L_2} \in \left[ 0; \frac{\pi}{2} \right] \\
    \cos \alpha = \left| \cos \varphi \right| = 
    \frac{\left| \overline{q_1} \cdot \overline{q_2} \right|}{\left| \overline{q_1} \right| \cdot \left| \overline{q_2} \right|}
\end{gather*}


\subsection{Взаимное расположение прямой линии $L$ и плоскости $P$}
\begin{gather*}
    L: \frac{x - x_0}{m} = \frac{y - y_0}{n} = \frac{z - z_0}{p} \\
    P: Ax + By + Cz + D = 0 \rightarrow \overline{N} = \{A, B, C\} \bot P
\end{gather*}

\subsubsection{Случай 1}
\[
    L \parallel P \Leftrightarrow \overline{q} \bot \overline{N}, 
    \overline{N} \cdot \overline{q} = 0 \\
    Am + Bn + Cp = 0
\]

\subsubsection{Случай 2}
\[
    L \nparallel P \Leftrightarrow Am + Bn + Cp \neq 0
\]
\begin{gather*}
    \beta = \widehat{L, P}: 
    \sin \beta = \left| \cos \varphi \right| = 
    \frac{\left| \overline{N} \cdot \overline{q} \right|}{\left| \overline{N} \right| \cdot \left| \overline{q} \right|}
\end{gather*}

\begin{gather*}
    k = L \cap P:
    \begin{cases}
        \frac{x - x_0}{m} = \frac{y - y_0}{n} = \frac{z - z_0}{p} = t \\
        Ax + By + Cz + D = 0
    \end{cases} \rightarrow
    \begin{cases}
        x = x_0 + mt \\
        y = y_0 + nt \\
        z = z_0 + pt
    \end{cases} \rightarrow
    \begin{cases}
        x_k = x_0 + mt_{\*} \\
        y_k = y_0 + nt_{\*} \\
        z_k = z_0 + pt_{\*}
    \end{cases} \\
    A(x_0 + mt) + B(y_0 + nt) + C(z_0 + pt) + D = 0 \rightarrow t = t_{\*}
\end{gather*}



\par\noindent\rule{\textwidth}{1pt}
\section{Практика}
\subsection{1.197(а)}

\paragraph{Дано}
\begin{gather*}
    L:
    \begin{cases}
        2x -  y + 2z - 3 = 0 \\
         x + 2y -  z - 1 = 0
    \end{cases} \overset{?}{\Rightarrow}
    \frac{x - x_0}{m} = \frac{y - y_0}{n} = \frac{z - z_0}{p}
\end{gather*}

\paragraph{Решение}
\paragraph{Пункт 1}
\begin{gather*}
    \overline{q} = 
    \overline{N_1} \times \overline{N_2} =
    \begin{vmatrix}
        \overline{i} & \overline{j} & \overline{k} \\
                  2  &          -1  &           2  \\
                  1  &           2  &          -1  
    \end{vmatrix} =
    \overline{i} \begin{vmatrix}
        -1 &  2 \\
        2  & -1
    \end{vmatrix} -
    \overline{j} \begin{vmatrix}
        2 &  2 \\
        1 & -1
    \end{vmatrix} +
    \overline{k} \begin{vmatrix}
        2 & -1 \\
        1 &  2
    \end{vmatrix} =
    -3 \overline{i} + 4 \overline{j} + 5 \overline{k} = \{ -3, 4, 5 \}
\end{gather*}

\paragraph{Пункт 2}
\begin{gather*}
    M_0: z = 0 \Rightarrow
    \begin{cases}
        2x -  y = 3 \\
         x + 2y = 1
    \end{cases} \\
    \text{Правило Крамера} \\
    \Delta = \det A = \begin{vmatrix}
        2 & -1 \\
        1 &  2  
    \end{vmatrix} = 4 + 1 = 5 \\
    \Delta_1 = \begin{vmatrix}
        3 & -1 \\
        1 &  2
    \end{vmatrix} = 7 \Rightarrow x = \frac{7}{5} =  1.4 \\
    \Delta_2 = \begin{vmatrix}
        2 & 3 \\
        1 & 1
    \end{vmatrix} -1 \Rightarrow y = \frac{-1}{5} = -0.2 \\
    M(1.4; -0.2; 0) \\
    L: \frac{x - 1.4}{-3} \frac{y + 0.2}{4} = \frac{z - 0}{5} \\
    \text{Проверка} \\
\end{gather*}



\subsection{1.200}
\paragraph{Дано}
\begin{gather*}
    L: \frac{x - 1}{2} = \frac{y}{1} = \frac{z + 1}{0} \\
    M_1(0; 1; 2) \\
\end{gather*}

\paragraph{Пункт 1 | Уравнение плоскости $P_1$}
\begin{gather*}
    P: \begin{cases}
        M_1 \in P \\
        M_0 \in P \\
        \overline{q} \parallel P
    \end{cases} \\
    \overline{M_0 M_1} = \{-1, 1, 3\} \\
    \begin{vmatrix}
        x - 1 & y - 0 & z + 1 \\
           -1 &     1 &     3 \\
            2 &     1 &     0 
    \end{vmatrix} = 0 \\
    (x - 1) \begin{vmatrix}
        1 & 3 \\
        1 & 0
    \end{vmatrix} -
    y \begin{vmatrix}
        -1 & 3 \\
         2 & 0
    \end{vmatrix} +
    (z + 1) \begin{vmatrix}
        -1 & 1 \\
         2 & 1 
    \end{vmatrix} = 0 \\
    -3x + 3 + 6y - 3z - 3 = 0 \\
    -3x + 6y - 3z = 0 \Rightarrow
    P_1: x - 2y + z = 0
\end{gather*}

\paragraph{Пункт 2 | Уравнение плоскости $P_2$}
\begin{gather*}
    P_2: M_1 \in P_2, P_2 \bot L \\
    \overline{N} = \overline{q} = \{2; 1; 0\} \\
    M_1 = \{0; 1; 2\} \\
    2(x - 0) + 1(y - 1) + 0(z - 2) = 0 \\
    P_2: 2x + y - 1 = 0
\end{gather*}

\paragraph{Пункт 3 | Уравнение прямой линии $L'$}
\begin{gather*}
    L': M_1 \in L', L' \bot L, L' \cap L \neq \varnothing \\
    L' = P_1 \cap P_2: \begin{cases}
         x - 2y + z = 0 \\
        2x +  y - 1 = 0
    \end{cases}
\end{gather*}

\paragraph{Пункт 4 | Расстояние от $M_1$ до $L$}
\begin{gather*}
    \rho(M_1, L) =
    \frac{\left| \overline{M_0 M_1} \times \overline{q} \right|}{\left| \overline{q} \right|} =
    \frac{\sqrt{3^2 + 6^2 + 3^2}}{\sqrt{2^2 + 1^2 + 0^2}} =
    \frac{\sqrt{54}}{\sqrt{5}} \\
    \overline{M_0 M_1} \times \overline{q} =
    \begin{vmatrix}
        \overline{i} & \overline{j} & \overline{k} \\
                  -1 &           1  &           3  \\
                   2 &           1  &           0
    \end{vmatrix} =
    \overline{i} \begin{vmatrix}
        1 & 3 \\
        1 & 0
    \end{vmatrix} -
    \overline{j} \begin{vmatrix}
        -1 & 3 \\
         2 & 0
    \end{vmatrix} +
    \overline{k} \begin{vmatrix}
        -1 & 1 \\
         2 & 1
    \end{vmatrix} =
    -3 \overline{i} + 6 \overline{j} - 3 \overline{k} = \{-3; 6; -3\}
\end{gather*}




%--------------------------------------------------
\newpage
\paragraph{ПЗ-13 | 2024.12.09 (3) (относительно конспекта Дани Палкина)}
\part{Кривые второго порядка}

\section{Теория}

\subsection{Общее уравнение}
\begin{gather*}
    Ax^2 + 2Bxy + Cyz + D + Ey + F = 0
\end{gather*}

\subsection{Приведение к каноническому виду}
\begin{gather*}
    B \neq 0 (x, y \text{есть})
\end{gather*}








\end{document}